\section{Related Work}
\label{sec:related}

\para{Verifiable instruction following.}
IFEval~\cite{zhou2023ifeval} introduced instructions whose compliance can be checked by code, such as ``use at least 300 words'' or ``do not use commas.''
InFoBench~\cite{qin2024infobench} decomposes instructions into checklists of requirements.
These benchmarks place one to three instructions in the user turn.
We borrow the idea of deterministic checkers but move the rules into the system prompt, scale their number to 40, and vary their surroundings.

\para{Many simultaneous instructions.}
FollowBench~\cite{jiang2024followbench} and ComplexBench~\cite{wen2024complexbench} add constraints step by step and find that performance falls as constraints accumulate.
\citet{harada2024curse} show that the rate of following {\em all} instructions falls roughly as the product of per-instruction success rates.
Our results agree that the all-rules-followed rate collapses quickly (to near zero by 15 rules for 8 of 9 models).
Unlike this line of work, we report {\em per-rule}, chance-corrected compliance, which reveals that each individual rule also becomes less likely to be followed as rules accumulate.
We also separate rule count from prompt length, which grow together when rules are simply appended.

\para{Position effects in long contexts.}
\citet{liu2024lost} show that models retrieve information best when it appears at the start or end of a long context, and worst in the middle.
It is natural to expect the same U-shape for instructions.
We find that this holds only partly: in long prompts the end of the prompt dominates, while in short numbered lists only the first rules gain.

\para{System prompts and instruction priority.}
\citet{wallace2024hierarchy} train models to prioritize system-level instructions over conflicting user instructions.
\citet{li2024instability} show that models drift away from system-prompt personas over long dialogs.
We study a complementary question: how well models follow many rules {\em within} a single system prompt, and what happens when two of those rules contradict each other.
